% ĐỀ KIỂM TRA CẤP 2 kiểu Sở HN: các bài baihn (Bài I, II, ...), ý nhỏ ychinh (1), 2), 3)). In đen trắng.
% Cách dùng (a): gọi thẳng mẫu rồi đặt môn bằng \mon.
% Đề có trắc nghiệm thì dùng thêm \phantn, tn... (xem Hướng dẫn soạn bài).
\documentclass[12pt]{article}
\usepackage{kiem_tra_thcs}
\begin{document}
\mon{TOÁN 9}
\so{Giỏi Toán cùng Trí Trí}
\truong{THCS CHƯƠNG DƯƠNG 25-26}
\kythi{Kiểm tra giữa học kì II}
\namhoc{2026 -- 2027}
\thoigian{90 phút}
\taodaude

\begin{baihn}{2}
Cho hai biểu thức $A=\dfrac{\sqrt{x}-3}{\sqrt{x}+2}$ và $B=\dfrac{5}{\sqrt{x}-3}$ với $x\ge0;\ x\ne9$.
\begin{ychinh}
\item Tính giá trị của biểu thức $A$ tại $x=25$.
\item Tính $P=A\cdot B$.
\end{ychinh}
\loigiai{
\begin{ychinh}
\item Thay $x=25$ vào $A$: $A=\dfrac{2}{7}$.
\item $P=\dfrac{5}{\sqrt{x}+2}$.
\end{ychinh}
}
\end{baihn}

\begin{baihn}{2,5}
\begin{ychinh}
\item Giải phương trình $x^2-5x+6=0$.
\item Một ô tô đi từ A đến B dài $120$ km hết $2$ giờ. Tính vận tốc ô tô.
\end{ychinh}
\loigiai{
\begin{ychinh}
\item $x=2$ hoặc $x=3$.
\item $v=120:2=60$ km/h.
\end{ychinh}
}
\end{baihn}

\hetde
\end{document}
